\textbf{结论名称：}椭圆焦点三角形面积公式 $S = b^{2}\tan\dfrac{\theta}{2}$

\textbf{适用条件：}椭圆标准方程 $\dfrac{x^{2}}{a^{2}}+\dfrac{y^{2}}{b^{2}}=1\;(a>b>0)$，$P$为椭圆上异于长轴端点的任意一点，$F_{1},F_{2}$为焦点，$\theta = \angle F_{1}PF_{2}$。

\textbf{变量说明：}$S$ 表示 $\triangle PF_{1}F_{2}$ 的面积；$b$ 为椭圆半短轴长；$\theta$ 为两焦半径 $PF_{1}$ 与 $PF_{2}$ 的夹角（$0<\theta<\pi$）。

\textbf{核心公式：}
\[
\boxed{S = b^{2}\tan\frac{\theta}{2}}
\]

\textbf{使用场景：}已知椭圆方程与两焦半径夹角，快速求出焦点三角形面积；或已知面积反推夹角；在选填题中避免联立方程组。

\textbf{简单例子：}椭圆 $\dfrac{x^{2}}{4}+\dfrac{y^{2}}{3}=1$，取 $P(0,\sqrt{3})$（短轴端点），则 $a=2,\;b=\sqrt{3},\;c=1$，$|PF_{1}|=|PF_{2}|=a=2$，焦距 $|F_{1}F_{2}|=2$。由余弦定理 $\cos\theta=\dfrac{2^{2}+2^{2}-2^{2}}{2\cdot2\cdot2}=\dfrac{1}{2}\Rightarrow\theta=\dfrac{\pi}{3}$。按公式 $S=b^{2}\tan\dfrac{\theta}{2}=3\tan\dfrac{\pi}{6}=3\cdot\dfrac{\sqrt{3}}{3}=\sqrt{3}$。直接计算 $S=\dfrac{1}{2}\cdot|F_{1}F_{2}|\cdot|y_{P}|=\dfrac{1}{2}\cdot2\cdot\sqrt{3}=\sqrt{3}$，结果一致。